% ============================================================
% CAPITOLO 2 -- INFRASTRUTTURA DI PROGETTO
% ============================================================
\section{Infrastruttura di progetto}

\subsection{Build e gestione delle dipendenze}

Il progetto Avro è organizzato come progetto multi-modulo Maven, con un
file \texttt{pom.xml} principale che coordina i sotto-progetti, tra cui
il modulo \texttt{lang/java/avro/}, dove abbiamo lavorato, che contiene
le classi \texttt{Resolver} e \texttt{Utf8} selezionate per il testing.

È stato necessario l'utilizzo di tre versioni diverse del JDK: 11, 17 e
21. Il progetto Avro richiede infatti la compatibilità con JDK 11 per il
codice compilato; inoltre, Randoop ed EvoSuite, generando codice sulla
base della JVM su cui vengono eseguiti, devono essere lanciati anch'essi
con JDK 11, per evitare di generare riferimenti a metodi non disponibili
nella versione richiesta per la compilazione finale. Le versioni JDK 17
e 21 sono invece necessarie per compilare ed eseguire il progetto nelle
altre fasi della build.

\subsection{Integrazione continua e misurazione della copertura}

È stata configurata un'integrazione continua tramite GitHub Actions: ad
ogni \texttt{push} sul branch principale, il file di configurazione
\texttt{isw2-ci.yml} esegue automaticamente una serie di step, tra cui
la compilazione del progetto e l'esecuzione dei test su una macchina
virtuale remota. Per la misurazione della copertura è stato utilizzato
JaCoCo, analizzando le metriche di statement e branch coverage, mentre
SonarCloud è stato integrato come supporto per l'analisi della qualità
del codice.

L'integrazione continua si è rivelata utile per individuare problemi non
visibili nell'ambiente di sviluppo locale: in particolare, un conflitto
tra il JDK usato per la compilazione in CI (JDK 21) e il codice generato
da EvoSuite (eseguibile correttamente solo su JDK 11) si è manifestato
soltanto durante l'esecuzione su GitHub Actions.

Sono stati inoltre utilizzati Spotless e Apache RAT per garantire la
qualità del codice e la conformità alle licenze. Spotless applica
automaticamente le regole di formattazione del progetto (indentazione,
ordine degli import), mentre Apache RAT verifica che ogni file sorgente
contenga l'header di licenza Apache 2.0, in linea con le politiche della
Apache Software Foundation. Entrambi i controlli vengono eseguiti prima
della fase di test, durante la compilazione: un file che non li rispetta
blocca l'intera build, indipendentemente dal fatto che i suoi test
vengano poi effettivamente eseguiti. Poiché i file di test generati
automaticamente da Randoop ed EvoSuite non rispettano lo stile del
progetto né includono l'header di licenza, è stato necessario
configurare esclusioni esplicite per queste cartelle in entrambi gli
strumenti.

\subsection{Isolamento dei test del progetto}

L'esecuzione dei test già presenti nel progetto è stata disabilitata,
mantenendo attiva una suite di test composta unicamente dai test
sviluppati per questo lavoro.